% Partie : méthodes
% Chapitre : induction mathématique
% Section : définitions inductives

\documentclass[../../../include/open-logic-section]{subfiles}

\begin{document}
\olfileid[fr]{mth}{ind}{idf}

\olsection{Définitions inductives}

En logique, nous définissons
très souvent des catégories
d'objets \emph{inductivement} :
nous donnons des règles qui
précisent ce qui compte
comme objet de la catégorie
et comment en former de
nouveaux à partir de ceux
que nous avons déjà.
Ainsi, certains types
de suites de symboles,
comme les termes et
les !!{formula}s d'un
langage, reçoivent
souvent une définition
inductive. Prenons
un exemple simple.
Considérons les suites
composées des lettres
$\mathrm{a}$, $\mathrm{b}$,
$\mathrm{c}$, $\mathrm{d}$,
du symbole~$\circ$
et des crochets $[$ et
$]$, telles que
«~$[[\mathrm{c} \circ \mathrm{d}][$~»,
«~$[\mathrm{a}[]\circ]$~»,
«~$\mathrm{a}$~» ou
«~$[[\mathrm{a} \circ \mathrm{b}]
\circ \mathrm{d}]$~».
Vous avez probablement
l'impression que les
deux premières ont
quelque chose d'incorrect :
dans la première,
les crochets ne
s'apparient pas ;
dans la seconde,
le symbole~$\circ$
ne relie pas deux
expressions qui ont
elles-mêmes un sens.
La troisième et la
quatrième paraissent
meilleures : chaque
crochet ouvrant $[$
possède un crochet
fermant $]$ correspondant,
s'il y a des crochets ;
et chaque $\circ$
relie deux expressions
correctes encadrées
par des crochets.

Nous voulons préciser
ce que sont ces
«~termes bien formés~».
Tout d'abord,
chaque lettre
prise seule
en est un.
Tout terme qui
n'est pas une
lettre seule
doit avoir la
forme
«~$[t \circ s]$~»,
où $s$ et $t$
sont eux-mêmes
bien formés.
Réciproquement,
si $t$ et $s$
le sont, nous
pouvons placer
$\circ$ entre
eux et entourer
le résultat
d'une paire
de crochets
pour former
un nouveau
terme bien
formé.
Ces opérations
permettent de
\emph{définir}
l'ensemble des
termes bien
formés : c'est
une \emph{définition
inductive}.

\begin{defn}[Termes bien formés]
  L'ensemble des
  \emph{termes bien formés}
  est défini inductivement
  par les règles suivantes :
  \begin{enumerate}
  \item Chaque lettre
    $\mathrm{a}$, $\mathrm{b}$,
    $\mathrm{c}$ ou $\mathrm{d}$
    est un terme bien formé.
  \item Si $s_1$ et $s_2$
    sont des termes bien
    formés, alors
    $[s_1 \circ s_2]$
    l'est aussi.
  \item Aucun autre objet
    n'est un terme
    bien formé.
  \end{enumerate}
\end{defn}

Selon cette définition,
une suite de symboles
est un terme bien
formé si et seulement
si elle peut être
construite par les
règles (1) et~(2)
en un nombre fini
d'étapes. À la
première étape,
nous construisons
les termes réduits
à une lettre :
\[
\mathrm{a}, \mathrm{b}, \mathrm{c}, \mathrm{d}.
\]
À la deuxième,
nous appliquons (2)
à ces termes :
\[
[\mathrm{a} \circ \mathrm{a}],
[\mathrm{a} \circ \mathrm{b}],
[\mathrm{b} \circ \mathrm{a}],
\dots, [\mathrm{d} \circ \mathrm{d}].
\]
Nous obtenons ainsi
toutes les combinaisons
de deux lettres.
À la troisième étape,
nous appliquons encore
(2) à deux termes
bien formés quelconques
déjà construits.
Nous obtenons, par
exemple,
$[\mathrm{a} \circ
  [\mathrm{a} \circ \mathrm{a}]]$,
où $t$ est le terme
$\mathrm{a}$ de la
première étape et
$s$ est le terme
$[\mathrm{a} \circ \mathrm{a}]$
de la deuxième.
Nous obtenons aussi
$[[\mathrm{b} \circ \mathrm{c}]
  \circ [\mathrm{d} \circ \mathrm{b}]]$
à partir des deux
termes
$[\mathrm{b} \circ \mathrm{c}]$
et
$[\mathrm{d} \circ \mathrm{b}]$
de la deuxième étape.
Et ainsi de suite.
La clause (3) empêche
qu'un objet non construit
de cette manière se
glisse dans l'ensemble.

Nous n'avons pas encore
démontré que toute suite
de symboles qui nous
paraît bien formée
l'est selon cette
définition. En revanche,
il devrait être clair
que chaque suite
effectivement construite
paraît bien formée :
ses crochets
s'apparient et
chaque $\circ$
relie des parties
elles-mêmes bien formées.

Une définition inductive
est particulièrement
utile quand on veut
démontrer une propriété
de tous les termes
bien formés : elle
indique quels cas
examiner. Si l'on
vous dit que $t$
est un tel terme,
la définition
inductive précise
ses formes possibles :
soit $t$ est une
lettre, soit
$t=[s_1 \circ s_2]$
pour deux termes
bien formés
$s_1$ et~$s_2$.
Grâce à la
clause (3),
il n'y a pas
d'autre possibilité.

La forme forte de
la récurrence devient
particulièrement utile
pour démontrer des
énoncés portant sur
tout un ensemble
défini inductivement.
Supposons, par exemple,
que nous voulions
montrer que, dans
chaque terme bien
formé de longueur~$n$,
le nombre de
crochets ouvrants
$[$ est $< n/2$.
On peut y voir
un énoncé portant
sur tous les~$n$ :
pour chaque~$n$,
tout terme bien
formé de longueur~$n$
contient moins de
$n/2$ crochets
ouvrants.

\begin{prop}
  Pour tout~$n$,
  le nombre de
  crochets ouvrants
  $[$ dans un
  terme bien formé
  de longueur~$n$
  est $< n/2$.
\end{prop}

\begin{proof}
Pour démontrer ce
résultat par récurrence
forte, il faut établir, pour tout entier naturel~$k$,
l'implication suivante :
\begin{quote}
  Si, pour chaque
  $l<k$, tout terme
  bien formé de
  longueur~$l$
  possède moins
  de $l/2$
  crochets ouvrants
  $[$, alors
  tout terme bien
  formé de
  longueur~$k$
  possède moins
  de $k/2$
  crochets
  ouvrants $[$.
\end{quote}
Supposons l'antécédent
de cette implication :
pour tout $l<k$,
les termes bien
formés de longueur~$l$
contiennent moins
de $l/2$ crochets
ouvrants. C'est
l'hypothèse de
récurrence. Montrons
que la même propriété
vaut pour les termes
de longueur~$k$.

Pour $k=0$, aucun terme bien formé n'a cette longueur :
l'énoncé est donc vérifié. Soit maintenant $t$ un terme bien
formé de longueur~$k$.
Sa définition inductive
ne laisse que deux
cas : (1)~$t$ est
une lettre seule ;
ou (2)~$t$ a
la forme
$[s_1 \circ s_2]$
pour deux termes
bien formés
$s_1$ et~$s_2$.
\begin{enumerate}
\item Si $t$ est
  une lettre,
  alors $k=1$
  et $t$ ne
  contient aucun
  crochet ouvrant.
  Comme $0<1/2$,
  la propriété
  est vérifiée.
\item Supposons
  $t=[s_1 \circ s_2]$
  pour deux termes
  bien formés
  $s_1$ et~$s_2$.
  Notons $l_1$
  la longueur
  de~$s_1$ et
  $l_2$ celle
  de~$s_2$.
  La longueur
  de~$t$ vaut
  $k=l_1+l_2+3$ :
  les longueurs
  des deux sous-termes,
  plus les trois
  symboles $[$,
  $\circ$ et $]$.
  Donc $l_1<k$
  et $l_2<k$.
  L'hypothèse de
  récurrence s'applique
  ainsi à $s_1$
  et à~$s_2$ :
  le nombre~$m_1$
  de crochets
  ouvrants dans~$s_1$
  est $<l_1/2$,
  et leur nombre~$m_2$
  dans~$s_2$
  est $<l_2/2$.

  Le nombre de
  crochets ouvrants
  de~$t$ est
  $m_1+m_2+1$ :
  ceux des deux
  sous-termes,
  plus le crochet
  ouvrant extérieur.
  Puisque
  $m_1<l_1/2$
  et $m_2<l_2/2$,
  nous obtenons :
  \[
  m_1 + m_2 + 1
  < \frac{l_1}{2}
    + \frac{l_2}{2} + 1
  = \frac{l_1+l_2+2}{2}
  < \frac{l_1+l_2+3}{2}
  = k/2.
  \]
\end{enumerate}
Dans chaque cas,
le nombre de
crochets ouvrants
de~$t$ est $<k/2$.
La récurrence
forte démontre
donc la proposition.
\end{proof}

\begin{prob}
  Définissez l'ensemble
  des termes
  «~super bien formés~»
  par les clauses
  suivantes :
  \begin{enumerate}
  \item Chaque lettre
    $\mathrm{a}$, $\mathrm{b}$,
    $\mathrm{c}$ ou $\mathrm{d}$
    est un terme
    super bien formé.
  \item Si $s$ est
    un terme super
    bien formé,
    alors $[s]$
    l'est aussi.
  \item Si $s_1$
    et $s_2$
    sont des termes
    super bien formés,
    alors
    $[s_1 \circ s_2]$
    l'est aussi.
  \item Aucun autre
    objet n'est un
    terme super
    bien formé.
  \end{enumerate}
  Montrez que le
  nombre de crochets
  ouvrants $[$
  dans un terme
  super bien formé~$t$
  de longueur~$n$
  est $\le n/2+1$.
\end{prob}

\end{document}
